% sheaf_method.tex — Non-constant sheaf formulation for TrianguLang
% This file can be swapped into method.tex Section 3.5 (Sheaf Consistency Loss)
% when ready. It extends the constant sheaf with learned restriction maps.

\subsection{Sheaf Consistency Loss}
\label{sec:sheaf}

To enforce cross-view prediction coherence, we introduce a \textbf{cross-view consistency loss} derived from the cellular sheaf Laplacian~\cite{currysheaf,sheafnn}. Unlike a standard pairwise consistency penalty, the sheaf framework provides a principled mechanism for \emph{view-dependent transformations} that account for geometric context at each correspondence.

\paragraph{Sheaf structure.} We model multi-view segmentation as a cellular sheaf $\mathcal{F}$ over a view graph $G = (V, E)$, where each vertex $v_i \in V$ corresponds to a camera view and edges $e_{ij} \in E$ connect views sharing geometric overlap. The \emph{stalks} are the data spaces assigned to each vertex and edge:
\begin{itemize}[nosep,leftmargin=*]
    \item \textbf{Vertex stalks:} $\mathcal{F}(v_i) = \mathbb{R}^d$, per-view feature vectors at each pixel (e.g., SAM3 encoder features, $d = 256$).
    \item \textbf{Edge stalks:} $\mathcal{F}(e_{ij}) = \mathbb{R}^{d_e}$, a shared comparison space for each view pair ($d_e = 32$).
\end{itemize}

\paragraph{Non-constant restriction maps.} The key departure from prior work is that our \emph{restriction maps} $\mathcal{F}_{v_i \trianglelefteq e_{ij}}: \mathcal{F}(v_i) \to \mathcal{F}(e_{ij})$ are \textbf{not identity}. Each restriction map consists of a shared linear projection modulated by a per-point geometric gate:
\begin{equation}
\mathcal{F}_{v_i \trianglelefteq e_{ij}}(\mathbf{f}_i, \mathbf{g}_i) = W \mathbf{f}_i \odot \sigma(\text{Gate}(\mathbf{g}_i))
\label{eq:restriction}
\end{equation}
where $W \in \mathbb{R}^{d_e \times d}$ is a shared projection matrix, $\mathbf{g}_i \in \mathbb{R}^{5}$ is the geometric context at the correspondence point (depth, correspondence distance, displacement, viewing angle proxy, depth edge indicator), and $\text{Gate}: \mathbb{R}^5 \to \mathbb{R}^{d_e}$ is a small MLP that outputs per-dimension gate values. The sigmoid $\sigma$ constrains gates to $(0, 1)$, and the elementwise product $\odot$ selectively attenuates feature dimensions based on geometric reliability.

\textbf{Why this is non-constant.} In a constant sheaf, restriction maps are identity, so every view contributes equally regardless of geometry. Eq.~\cref{eq:restriction} defines a \emph{non-constant} sheaf because: (1) the gate varies per spatial location based on the geometric context $\mathbf{g}_i$, and (2) different views at the same 3D point have different geometric contexts (depth, viewing angle, edge proximity), yielding different effective projections into the edge space. This enables the sheaf to learn that predictions near depth discontinuities or at grazing viewing angles are less reliable and should be down-weighted in the consistency comparison.

\paragraph{Initialization.} The gate bias is initialized to $+2.0$, so $\sigma(\text{Gate}(\mathbf{0})) \approx 0.88 \approx 1$. At the start of training, the restriction maps approximate a constant sheaf (pure linear projection), and the model learns deviations from this baseline as needed.

\paragraph{Sheaf Laplacian energy.} The consistency loss is the sheaf Laplacian quadratic form:
\begin{equation}
\mathcal{L}_{\text{sheaf}} = \frac{1}{|E|} \sum_{e_{ij} \in E} \frac{1}{|\mathcal{C}_{ij}|} \sum_{(\mathbf{u}_i, \mathbf{u}_j) \in \mathcal{C}_{ij}} \left\| \mathcal{F}_{v_i \trianglelefteq e_{ij}}(\mathbf{f}_i^{\mathbf{u}}, \mathbf{g}_i^{\mathbf{u}}) - \mathcal{F}_{v_j \trianglelefteq e_{ij}}(\mathbf{f}_j^{\mathbf{u}}, \mathbf{g}_j^{\mathbf{u}}) \right\|_2^2
\label{eq:sheaf_feature}
\end{equation}
where $\mathcal{C}_{ij}$ are mutual nearest-neighbor pixel correspondences with 3D distance $< 10$cm, established via DA3-estimated pointmaps. When $\mathcal{L}_{\text{sheaf}} = 0$, the per-view predictions form a \emph{global section} of the sheaf: the multi-view system is perfectly consistent.

\paragraph{Connection to the sheaf Laplacian.} Eq.~\cref{eq:sheaf_feature} is precisely the quadratic form $\mathbf{s}^\top L_\mathcal{F} \mathbf{s}$ associated with the sheaf Laplacian $L_\mathcal{F} = \delta^\top \delta$, where $\delta: C^0(\mathcal{F}) \to C^1(\mathcal{F})$ is the coboundary operator mapping vertex sections to edge coboundaries via the restriction maps. The kernel $\ker(L_\mathcal{F})$ is the space of global sections $H^0(G; \mathcal{F})$. For a constant sheaf on a connected graph, $\dim(H^0) = 1$ (all views must agree exactly). For our non-constant sheaf, the learned gates expand $\ker(L_\mathcal{F})$ by down-weighting unreliable correspondences, allowing the model to tolerate geometric noise without penalizing correct predictions.

\paragraph{Spectral interpretation.} The eigenvalues of $L_\mathcal{F}$ characterize the sheaf's consistency structure. We verify empirically that: (1) $L_\mathcal{F}$ is positive semidefinite (guaranteed by construction as $\delta^\top \delta$), (2) the constant sheaf and learned sheaf produce different spectral profiles, confirming that training learns genuinely non-trivial restriction maps, and (3) the smallest non-zero eigenvalue (spectral gap) increases during training, indicating that the sheaf converges toward a structure where consistent predictions are well-separated from inconsistent ones.

\begin{remark}[Complementary consistency mechanisms]
GASA (Eq.~\cref{eq:gasa}) and the sheaf loss (Eq.~\cref{eq:sheaf_feature}) enforce consistency through complementary mechanisms. GASA's geometric kernel operates on \emph{features during inference}, gating cross-view attention based on 3D distance. The sheaf loss operates on \emph{feature correspondences during training}, directly penalizing inconsistency at known 3D matches through learned, geometry-aware restriction maps. The two mechanisms are complementary: GASA provides soft geometric gating that works without explicit correspondences, while the sheaf loss provides exact correspondence-based supervision that teaches the encoder to produce geometrically consistent features.
\end{remark}
